\section{The graphs of Horak and Kim}\label{sec:examples}

The lower bound~\eqref{eq:hk} is due to Horak and Kim~\cite[Theorem~5.1]{HK}. We recall their
graphs for odd rank, because they show where the second term comes from.

Let $j\ge2$. Take vertices $a_i,b_i,c_i,d_i$ for $i$ modulo $j$. Join each of $a_i,b_i$ to each of
$c_i,d_i$; these four edges are the \emph{square} $i$. Join $c_i$ to $a_{i+1}$ and $d_i$ to
$b_{i+1}$; these two edges are the \emph{interface} $i$. The graph $M_j$ is cubic, with $4j$
vertices, $6j$ edges and rank $r=2j+1$. In the notation of~\cite{HK} it is the graph $MC_{2j}$.

\begin{proposition}[Horak and Kim]\label{prop:hk}
For $j\ge2$,
\[
  c(M_j)\;=\;2^{2j}+\bigl(2j+\tfrac12\bigr)2^{j}-3j\;=\;2^{r-1}+\bigl(r-\tfrac12\bigr)2^{(r-1)/2}-\tfrac32(r-1).
\]
\end{proposition}

\begin{proof}
The two edges of an interface form a cut together with the two edges of any other interface. So a
cycle uses an odd number of edges of every interface, or an even number of edges of every
interface.

\emph{One edge of every interface.} There are $2^{j}$ ways to choose the edges. In each square the
cycle then runs from one of $a_i,b_i$ to one of $c_i,d_i$, and there are two such paths: the edge,
and the path through the other two vertices of the square. Every choice gives a cycle. This gives
$2^{j}\cdot2^{j}=2^{r-1}$ cycles.

\emph{An even number of edges of every interface.} If the cycle uses no interface, it is one of the
$j$ squares. If it uses both edges of all interfaces, it uses two disjoint edges of every square,
and there are two ways to choose them in each square. Such a choice is one cycle, and not two, for
half of the $2^{j}$ choices: $2^{j-1}$ cycles. Otherwise the interfaces that it uses are $d$
consecutive ones, with $1\le d\le j-1$. There are $j$ such sets for each $d$. The cycle closes in
the two squares at the ends, in two ways each, and it uses two disjoint edges in each of the
$d-1$ squares between them. This gives $j\,2^{d+1}$ cycles for each $d$. Together,
\[
  j+2^{j-1}+j\sum_{d=1}^{j-1}2^{d+1}\;=\;\bigl(2j+\tfrac12\bigr)2^{j}-3j .\qedhere
\]
\end{proof}

The first class contains exactly half of the even subgraphs of $M_j$, and all of them are cycles.
The second term counts the cycles in the other half. Most of these are the cycles that use a
proper set of consecutive interfaces. They exist because each square is a cycle with a cut of four
edges around it. This is the configuration that Theorem~\ref{thm:smallcut} treats.

The graphs $M_j$ also show that the conclusion of Theorem~\ref{thm:half} fails without the
hypothesis on the girth. Let $W$ be a square of $M_j$ and let $e$ be one of its edges. The cycles
of the first class through $e$ number $2^{r-2}$, and the square itself is a further cycle through
$e$. So $c_e(M_j)>2^{r-2}$, and Lemma~\ref{lem:edge} gives $q(M_j,V(W))>\half$. Here $M_j-V(W)$
contains cycles of length four.

For even rank, Horak and Kim use the same construction on two cycles of odd length, with one
uncrossed pair of rungs~\cite[Section~5]{HK}. We tested both formulas of~\eqref{eq:hk} by computer
for small ranks (Appendix~\ref{app:checks}).
